\documentclass[11pt,a4paper]{article}

\usepackage[T1]{fontenc}
\usepackage{lmodern}
\usepackage[margin=2.5cm]{geometry}
\usepackage{amsmath,amssymb}
\usepackage{booktabs}
\usepackage{tabularx}
\usepackage{listings}
\usepackage{xcolor}
\usepackage[hidelinks]{hyperref}

\lstset{
  basicstyle=\ttfamily\small,
  columns=fullflexible,
  keepspaces=true,
  frame=single,
  language=Java
}

\newcommand{\code}[1]{\texttt{#1}}

\title{Logic-Based Deletion Analysis of Java Code Dependencies\\[0.3em]
\large A Prototype for Giving Feedback to AI Coding Assistants}
\author{Mathis Burger}
\date{24 October 2026}

\begin{document}

\maketitle

\begin{abstract}
AI coding assistants change and delete code without knowing reliably what
depends on it. This document describes a prototype that encodes code-level
dependencies of a Java project as a feature model, translates it to
propositional logic and answers deletion questions on this logical
encoding: \emph{Can this class or member be deleted? If not, which
functionality breaks, and why?} The encoding and five queries (including
two MaxSAT queries and an explanation via a minimal unsatisfiable subset)
are specified and implemented on top of the SAT solver Sat4j. For the
constraints implemented so far, the encoding lies in a fragment that
propagation alone can solve; where a SAT solver becomes necessary is one
of the open research questions. A source extractor based on
JavaParser builds the feature tree and the plain references between code
elements. The queries have been checked on a small hand-built example.
An end-to-end run from source on the same example works, but answers only
some of the queries correctly, because the extractor does not yet produce
roots and polymorphism constraints. A run on a real project and an
evaluation have not been done.
\end{abstract}

\begin{center}
\fbox{\begin{minipage}{0.92\textwidth}
\textbf{Nature of this document.}
This is a write-up of a prototype, not a research paper. Its purpose is
to show that the research I propose is implementable and to make its
central idea concrete. It does not yet show that the approach works on
real projects: no evaluation has been run, and the extractor covers only
part of the encoding. The related work in Section~\ref{sec:related} is
preliminary and mostly based on abstracts.

\medskip
\textbf{Use of AI.}
AI assistants (Claude Code) were used to draft text and code in this
project. I checked all code and text personally. The research idea and the
design decisions are my own.
\end{minipage}}
\end{center}

\section{Motivation and Question}

When an AI coding assistant proposes to delete a class, it usually learns
whether this was safe only after compiling and running the tests, or not
at all. A tool that answers deletion questions \emph{before} the change,
and explains its answer, could serve as feedback for the assistant's
decision.

The question behind the prototype is:

\begin{quote}
Can a propositional encoding of code-level dependencies, solved with a
SAT solver where necessary, tell an AI coding assistant, before
compilation or testing, whether a proposed deletion breaks the project?
Concretely: can class or member $C$ be deleted? If not, which
functionality breaks, and why?
\end{quote}

The phrase ``where necessary'' matters: as Section~\ref{sec:complexity}
shows, the constraints implemented so far do not require a SAT solver.
Finding out which rules about code do require one is part of the
question.

The approach borrows from two areas. Package managers have long encoded
dependencies between packages as propositional formulas and solved them
with SAT solvers. Software product lines describe valid combinations of
features as feature models, which also translate to propositional logic.
The prototype applies both ideas at the level of classes, fields,
constructors and methods.

\section{Preliminary Related Work}\label{sec:related}

This section lists the work found so far. It is not a systematic review.
For most entries only the abstract has been read.

\paragraph{Feature models and propositional logic.}
Batory~\cite{batory2005} connects feature models, grammars and
propositional formulas, so that feature models can be translated to
propositional logic and analysed with off-the-shelf SAT solvers. Kuiter et
al.~\cite{kuiter2022} show that the choice of CNF transformation affects
the correctness and performance of feature-model analyses; the encoding
here avoids the question because every constraint is already a single
clause. Sundermann et al.~\cite{sundermann2023} evaluate \#SAT solvers on
industrial feature models, which becomes relevant if one wants to count
valid deletion sets.

\paragraph{SAT on code-level references.}
The closest prior work found is \emph{safe composition} by Thaker et
al.~\cite{thaker2007}: feature modules of a product line reference classes,
methods and fields defined in other modules, and SAT is used to check that
no product allowed by the feature model contains a reference to an
undefined element. Apel et al.~\cite{apel2010} and K\"astner et
al.~\cite{kaestner2012} give type systems that guarantee all variants of a
product line are well-typed. How the proposed research relates to this
line of work is discussed in Section~\ref{sec:questions}.

\paragraph{SAT-based package dependency solving.}
Mancinelli et al.~\cite{mancinelli2006} encode package installability as a
SAT problem, with one Boolean variable per package version, and check
whole Debian and Mandriva repositories with it. Abate et
al.~\cite{abate2020} survey dependency solving in package managers, note
that it is NP-complete in non-trivial component models, and observe that
SAT-based solving is spreading. Sat4j~\cite{leberre2010}, the solver used
here, provides SAT, MaxSAT and pseudo-Boolean solving in Java and has been
used for Eclipse's plugin dependency management. Pinckney et
al.~\cite{pinckney2023} use Max-SMT to optimise npm dependency solving.

\paragraph{Breaking changes and LLM feedback.}
BUMP~\cite{reyes2024} is a benchmark of reproducible breaking dependency
updates in Java projects. Roseau
\cite{latappy2025} detects API breaking changes in Java libraries.
CodePlan~\cite{bairi2023} uses static dependency analysis to plan
repository-wide edits made by an LLM; it does not use SAT.

\section{Approach}

The full specification is in \code{docs/encoding.md} in the repository.
This section summarises it.

\subsection{Features}

Every feature is one Boolean variable: \emph{true} means the element is
present in the code, \emph{false} that it is deleted. The feature tree is
\[
  \text{project root} \to \text{package} \to \text{class}
  \to \text{field / constructor / method},
\]
with named nested classes as children of their enclosing class. Each
overload of a method or constructor is its own feature, identified by its
parameter types (e.g.\ \code{Report\#total(Shape)}). Parameters are not
separate features. Anonymous classes, local classes and lambdas are not
features; their references are attributed to the enclosing member. JDK and
library types are not modelled.

\subsection{Constraints}

Every constraint is a single clause; $a \to b$ stands for
$\neg a \lor b$.
\begin{itemize}
  \item \textbf{Tree:} the project root is present; every child requires its
        parent ($\mathit{child} \to \mathit{parent}$). All other features are
        optional and therefore deletion candidates.
  \item \textbf{References:} the innermost feature containing a reference
        (its \emph{source}) requires the referenced feature: method and
        constructor calls (bound to the overload the compiler resolves),
        type usages, inheritance, annotations and field accesses.
  \item \textbf{Polymorphism:} for an interface or abstract class $T$ with
        implementations $C_1,\dots,C_n$, an or-group
        $T \to (C_1 \lor \dots \lor C_n)$; a call on an abstract method
        requires the declaration and at least one implementation; a concrete
        class requires its implementation of every inherited abstract method.
        Lambdas count as implementations through their enclosing member.
        Groups without any implementation in the project are skipped.
  \item \textbf{Roots} (``functionality''): every \code{main} method and
        every test class is a root and must be present; a root class
        requires all of its members.
  \item \textbf{Deletion} of a set $D$: unit clauses $\neg d$ for
        $d \in D$.
\end{itemize}

\subsection{Queries}

Let $F$ be all clauses except the root unit clauses, $R$ the root unit
clauses and $D$ the deletion set.

\begin{center}
\begin{tabularx}{\textwidth}{@{}l X X@{}}
\toprule
& Question & Query \\
\midrule
Q1 & Can $D$ be deleted alone? &
  SAT of $F \land R \land \neg D$ under the assumption that every feature
  outside the deleted subtrees stays present \\
Q2 & Can $D$ be deleted together with dependent code? &
  SAT of $F \land R \land \neg D$ \\
Q3 & Which code must go with $D$? &
  MaxSAT: $F \land \neg D$ hard, one soft unit clause per remaining feature \\
Q4 & Which functionality breaks? &
  MaxSAT: $F \land \neg D$ hard, one soft unit clause per root \\
Q5 & Why? &
  Minimal unsatisfiable subset (MUS) of the Q2 formula; each clause keeps
  the code reference that produced it \\
\bottomrule
\end{tabularx}
\end{center}

\subsection{Complexity, and why SAT is used anyway}\label{sec:complexity}

Every clause of the encoding has at most one negative literal, so the
formulas are \emph{dual-Horn}. Their satisfiability can be decided in
linear time: negating every variable turns them into Horn formulas, for
which Dowling and Gallier~\cite{dowling1984} give a linear-time
algorithm.
Moreover, the set of models is closed under union: if $a$ is true in
$M_1 \cup M_2$, it is true in some $M_i$, which then makes one of the
positive literals of every clause $\neg a \lor b_1 \lor \dots \lor b_k$
true in $M_i$ and therefore in the union. Hence a unique maximal model
exists, computable by propagating ``absent'' from $D$ to a fixpoint, and it
answers Q2--Q4 at once. This is only a proof sketch and still has to be
checked carefully.

For the current encoding, a graph or propagation algorithm therefore gives
the same answers as SAT and MaxSAT. Explanations are not an argument for
SAT either: in this fragment, the chain of implications that forced a
feature to be absent can be read directly from the propagation. Using a
solver library is merely convenient here, because it provides minimal
unsatisfiable subsets without extra implementation effort. The actual
reason to keep a SAT solver is that planned user-defined rules
(``exactly one implementation'', ``A excludes B'') contain clauses with
more than one negative literal and therefore leave the dual-Horn
fragment, so that efficient propagation is no longer guaranteed. This
reason is prospective: whether such rules are needed in practice has not
been shown. The propagation algorithm will serve as a baseline in the
evaluation.

\section{Prototype}

The prototype is a Maven project in Java~21 with Sat4j~2.3.6 and
JavaParser~3.28.2 (including its symbol solver), tested with JUnit~5.

\begin{itemize}
  \item \code{FeatureModel} and \code{Constraint} hold the model in clause
        form. Each clause keeps a human-readable reason, so that a MUS can be
        traced back to the code.
  \item \code{DeletionQueries} implements Q1--Q5: Q1 and Q2 with Sat4j's
        default solver (Q1 via assumptions), Q3 and Q4 with Sat4j's weighted
        MaxSAT decorator, Q5 with Sat4j's \code{Xplain}.
  \item \code{Extractor} builds the model from source in layers. Implemented
        are layer~1 (the feature tree) and layer~2 (inheritance, type usage,
        annotations, field access, constructor calls, resolved with the
        JavaParser symbol solver). Names that cannot be resolved, such as
        library annotations, are reported as warnings.
\end{itemize}

Not implemented are layer~3 (method calls and implicit \code{super()}
calls), layer~4 (polymorphism) and layer~5 (roots); see
Section~\ref{sec:limitations}. The pipeline nevertheless runs end to end
on the partial model; Section~\ref{sec:checked} shows what it gets right
and wrong.

\section{Worked Example}

Example A is a small synthetic program with an interface and two
implementations:

\begin{lstlisting}
interface Shape { double area(); }
class Circle implements Shape { public double area() { return 1; } }
class Square implements Shape { public double area() { return 1; } }
class Report { static double total(Shape s) { return s.area(); } }
class ReportTest { @Test void total() { Report.total(new Circle()); } }
\end{lstlisting}

Its model has 12 features (project root, package \code{p}, the five classes
and their five members) and 24 clauses. The root is \code{ReportTest}.
Three deletions were worked out by hand:

\begin{center}
\begin{tabular}{@{}lllll@{}}
\toprule
Deletion & Q1 & Q2 & Q3 cascade & Q4 broken roots \\
\midrule
\code{Square} & yes & yes & \code{Square}, \code{Square\#area()} & none \\
\code{Circle} & no & no & \code{Circle}, \code{Circle\#area()}, & \code{ReportTest} \\
 & & & \code{ReportTest\#total()}, \code{ReportTest} & \\
\code{Square\#area()} & no & yes & \code{Square\#area()}, \code{Square} & none \\
\bottomrule
\end{tabular}
\end{center}

\begin{itemize}
  \item \code{Square} can be deleted because \code{Circle} still satisfies
        the or-group of \code{Shape} and the call \code{s.area()}. This is
        the case the or-groups exist for.
  \item \code{Circle} cannot be deleted. The explanation (Q5) is the MUS
        \{\code{ReportTest} is a root; \code{ReportTest} requires
        \code{ReportTest\#total()}; \code{ReportTest\#total()} calls
        \code{new Circle()}; \code{Circle} is deleted\}.
  \item \code{Square\#area()} cannot be deleted alone, because
        \code{Square} would no longer implement \code{Shape}; it can be
        deleted together with \code{Square}.
\end{itemize}

\section{What Has Been Checked}\label{sec:checked}

All numbers in this section come from \code{mvn test} in the repository
(12 tests, all passing on 24 October 2026).

\begin{itemize}
  \item \code{ExampleATest} (4 tests): the hand-built model of Example A
        consists of exactly the 24 documented clauses, and Sat4j's answers to
        Q1--Q5 for the three deletions above equal the hand-derived ones,
        including the exact MUS.
  \item \code{ExtractorTreeTest} (3 tests): the extracted feature tree of
        Example A equals the hand-built one; a second fixture covers
        overloads, generic and varargs parameters, nested classes, enums,
        interfaces, records and nested packages, and checks that anonymous
        and local classes, lambdas and initializer blocks do not become
        features; variable numbering is deterministic.
  \item \code{ExtractorReferenceTest} (2 tests): layer~2 produces exactly the
        four plain-reference clauses of Example A, and 15 hand-derived
        clauses for a fixture covering inheritance, annotations, generic
        field types, enum constants, static initializers, explicit
        \code{super(\dots)}, implicit and explicit constructors, casts,
        class literals, inherited fields, and references from lambdas and
        anonymous classes.
  \item \code{EndToEndTest} (3 tests): the sources of Example A are
        extracted with layers~1 and~2 only, and the queries run on the
        resulting model. The test asserts the answers derived by hand for
        this partial model (\code{docs/encoding.md}, Section~5), including
        the ones that are wrong compared with the full model.
\end{itemize}

These tests show that the encoding, the queries and the first two
extractor layers behave as specified on small examples.

\paragraph{End to end with layers 1 and 2 only.}
The extracted model of Example A contains the feature tree and the four
plain references, but no method calls, no or-group, no implementation
obligations and no roots. Its answers compare with those of the full model
as follows:

\begin{center}
\begin{tabular}{@{}llll@{}}
\toprule
Deletion & Q1 alone & Q2 with cascade & Same as full model \\
\midrule
\code{Square} & yes & yes & yes, Q1--Q4 \\
\code{Circle} & no & yes (full: no) & Q1 only \\
\code{Square\#area()} & yes (full: no) & yes & Q2 and Q4 only \\
\bottomrule
\end{tabular}
\end{center}

Deleting \code{Circle} is still correctly rejected by Q1, because the test
calls \code{new Circle()}. Without roots, however, nothing has to stay
present: Q2 accepts the deletion, the cascade (Q3) misses
\code{ReportTest}, no functionality is reported broken (Q4), and there is
nothing to explain (Q5). Without the implementation obligation, even Q1
accepts deleting \code{Square\#area()}, although \code{Square} would then no
longer compile. The run therefore shows that the pipeline works from
source to answer, and which of the missing layers each answer depends on:
roots (layer~5) for Q2, Q4 and Q5, and polymorphism (layer~4) even for
Q1.

\section{Open Questions and How the Research Would Address Them}\label{sec:questions}

The prototype raises several questions that it cannot answer itself. This
section states them openly, together with how I would address them in
the research. The points in it are arguments and hypotheses, not results.

\paragraph{Why not simply compile and run the tests?}
Compiling and testing remain the ground truth, and the evaluation would
use them as such. The tool is meant to complement them, not to replace
them. A compiler answers a question only after the edit has been made,
and it reports the errors of that one edit. The queries in this
prototype answer questions the compiler does not answer directly: whether
a deletion is possible \emph{before} the code is touched, what else would
have to be removed with it (Q3), which tests and entry points would be
lost (Q4), and why (Q5), as a chain of references rather than a list of
error messages. A single model can also be asked about many candidate
deletions without rebuilding the project each time. Whether this is
faster than incremental compilation in practice, and whether the
explanations help an AI assistant, are open questions that the
evaluation would have to measure.

\paragraph{Why SAT, if the current encoding is dual-Horn?}
For the constraints implemented so far, a SAT solver is not necessary
(Section~\ref{sec:complexity}), and the explanations could also be
obtained from propagation. I do not see this as a reason to drop the
approach, but it changes how the research question should be read: the
contribution is a \emph{logical encoding} of code-level dependencies
together with deletion queries and explanations, with SAT as a general
back end for the cases that need it. A central part of the research would be to find out where the boundary
lies, that is, which kinds of rules about code keep the problem in an
efficiently solvable fragment and which do not. Candidates that leave the
dual-Horn fragment are user-defined rules such as ``exactly one
implementation of this interface'' or ``these two components exclude each
other''. Whether such rules occur often enough in practice to matter is
itself an open question. If they do not, a propagation algorithm would be
sufficient, and that would be a result worth reporting as well. Until
this is known, I favour SAT as the back end, because rules outside the
dual-Horn fragment can then be added as further clauses without changing
the solver. Such constraints (e.g.\ $A \land B \to \neg C$) are better
handled by SAT than by a custom graph algorithm: propagation alone no
longer finds the answer, so the graph algorithm would have to add the
search that a SAT solver already provides. How well the SAT and especially the MaxSAT queries scale on
real projects has not been measured yet.

\paragraph{What is new compared to safe composition?}
Thaker et al.~\cite{thaker2007} already encode code references as
propositional constraints between feature modules and check them with
SAT. The research would build on
this idea rather than claim it. The first intended difference is the
granularity: in safe composition a variable stands for a feature module,
and a reference only yields an implication between modules, whereas here
every class, field, constructor and method is a variable of its own. This
is what makes the deletion of a single element expressible at all.
Further differences are the input (an ordinary project from which the model is
extracted automatically, not a product line with a hand-written feature
model), the question (the impact of one concrete deletion, the cascade
and the minimal set of broken functionality, rather than type safety of
all products) and the consumer (an AI assistant that needs an answer and
an explanation). Whether these differences hold up can only be judged
after a thorough reading of this line of work, which I plan as a first
step.

\paragraph{Where is the AI?}
The prototype contains no AI. It focuses on the part that has to work
first: a reliable answer to a deletion question. The AI-related question,
whether an assistant makes better decisions when it can ask the tool,
is a separate experiment. It would only make sense once the tool itself
has been evaluated against compilation and test results.

\paragraph{Does it work on real code?}
Only for one five-class example, and only partly
(Section~\ref{sec:checked}). The first milestone of the research would be
a complete end-to-end run: completing the extractor, deleting each class of a small
open-source subject in turn, and comparing the prediction with the result
of compiling and running its tests. False positives and false negatives
would be reported as they are, including the ones the current design
already predicts (Section~\ref{sec:limitations}).

As a first subject I propose
\href{https://github.com/iluwatar/java-design-patterns}{java-design-patterns},
an open-source collection of the common Java design patterns. Its modules
are small and well understood, so answers can be checked by hand, and the
patterns they demonstrate recur in real Java code. The modules themselves,
however, are demonstrations and are not structured like real projects. To
cover that, a large open-source project such as Apache Kafka would be
analysed as well; judging whether an answer is correct there requires
in-depth knowledge of the codebase, which would have to be built up
first. A more focused variant of the research would
first analyse which design patterns occur in practice and how relevant
they are, and then study how each of them is represented in the
encoding.

Two points of the evaluation design need care. First, the ground truth
differs per query. For Q1, it is the result of deleting only $C$ and then
compiling and running the tests. For Q2 and Q3, it is the result of
deleting the predicted cascade instead: the prediction is correct if the
project then compiles and only the predicted roots (Q4) are lost.
Second, the answers are likely to be imbalanced. In the
\code{strategy} module of the planned subject, for example, the
\code{main} method uses every strategy class, so almost none of them can
be deleted. A trivial
predictor that always answers ``not deletable'' would then look
accurate. The evaluation would therefore report measures that account
for imbalance, such as precision and recall per answer, rather than
accuracy alone, and compare against three baselines of increasing
strength: ``always not deletable''; plain reachability, which treats every
reference as required and ignores or-groups, so it cannot see that
another implementation takes over; and the propagation algorithm of
Section~\ref{sec:complexity}, which gives the same answers as the
encoding in the dual-Horn fragment and therefore isolates the cost and
benefit of using a solver. Subjects with more deletable code would also
be needed.

\paragraph{Is deletion too narrow?}
Deletion is a deliberately small starting point because its ground truth
is clear. Other changes can partly be expressed through it: changing a
method's signature, for example, amounts to deleting one overload and
adding another. Whether the approach extends to other kinds of edits is
part of the future work.

\section{Limitations}\label{sec:limitations}

\begin{itemize}
  \item \textbf{Partial end-to-end result, no evaluation.} The extractor
        lacks method calls, polymorphism and roots. From source, only one
        synthetic example has been run, and several of its answers are wrong
        for that reason (Section~\ref{sec:checked}). The planned subject
        (java-design-patterns, Section~\ref{sec:questions}) has not been
        run.
  \item \textbf{Resolved overload only.} A call requires exactly the
        overload the compiler picks today. If \code{bar(int)} is deleted,
        \code{bar(5)} still compiles against \code{bar(long)}, but the model
        reports a break (false positive).
  \item \textbf{Implicit \code{super()}} calls are not modelled yet, so
        deleting a superclass's no-argument constructor can break
        compilation unnoticed (false negative).
  \item \textbf{Static analysis only.} Reflection, dependency injection,
        \code{ServiceLoader}, references from resource files, generated code
        (e.g.\ Lombok), runtime-created implementations (mocks, proxies) and
        breakage that compiles but fails a test are out of scope. Mocks in
        particular cause false positives: a test that uses only a Mockito
        mock of an interface still passes when every implementation is
        deleted, but the or-group reports it as broken.
  \item \textbf{Approximate signatures.} Overloads are identified by
        parameter types as written in the source, so type variables are not
        erased to their bounds and equally named types are not told apart.
        Identifying each parameter type by its resolved, fully qualified
        name would fix this.
  \item \textbf{Dual-Horn fragment.} As shown above, the current encoding
        does not need a SAT solver; the benefit of SAT has to be shown with
        rules outside this fragment (Section~\ref{sec:questions}).
  \item \textbf{Preliminary literature review.} The related work is
        incomplete and mostly based on abstracts.
\end{itemize}

\section{Next Steps for the Research}

\textbf{NOTE:} Order is not prioritised.

\begin{itemize}
  \item Complete the extractor (method calls, implicit \code{super()},
        polymorphism, roots) and run the queries end to end.
  \item Implement the propagation baseline and compare it with SAT.
  \item Evaluate on open-source projects, starting with
        java-design-patterns: delete each class in turn, compare
        the prediction with compiling and running the tests (for Q2 and Q3:
        after deleting the predicted cascade), and report precision and
        recall per answer, false positives and negatives, and solver time,
        compared against trivial baselines and the propagation algorithm.
  \item Add user-defined rules in UVL constraint syntax (e.g.\ exactly one
        implementation), which leave the dual-Horn fragment, and determine
        the complexity of the resulting problem.
  \item Render explanations as text for an AI assistant and study whether
        the assistant decides better with the tool.
  \item Extend from single projects to services in several languages. Many
        dependencies between services (e.g.\ HTTP calls built from
        configuration values) cannot be proven by static analysis. An AI
        could infer them, and the encoding would need to include them
        together with their uncertainty, so that an answer shows which parts
        rest on proven dependencies and which on inferred ones.
  \item Carry out a proper literature review, starting with reading the
        safe-composition and variability-aware type-checking work in full.
\end{itemize}

\bibliographystyle{plainurl}
\bibliography{references}

\end{document}
